\documentclass[conference]{IEEEtran}
\usepackage{cite}
\usepackage{amsmath,amssymb,amsfonts}
\usepackage{array}
\usepackage{booktabs}
\usepackage{graphicx}
\usepackage{xcolor}
\usepackage{url}
\usepackage{multirow}
\newcommand{\origin}{\textsc{origin}}
\newcommand{\designonly}{\textsc{design-only}}
\begin{document}

\title{Where Do Human Emotions Originate?\\
A Distributed Brain-Body Research Program}

\author{\IEEEauthorblockN{Anonymous Research Artifact}
\IEEEauthorblockA{\textit{Affective Neuroscience Evidence Synthesis}\\
IEEE-style research report\\
2026}}

\maketitle

\begin{abstract}
The question of where human emotions originate is often posed as a search for an emotion center in the brain. That framing is not supported by contemporary affective neuroscience. This evidence synthesis and prospective research program instead treats emotion episodes as distributed brain-body control processes. Conserved neural systems support defense, reward, pain avoidance, arousal, autonomic regulation, and action selection. Cortical, hippocampal, and interoceptive systems integrate sensory context, bodily state, memory, learning, value, and behavioral goals. The resulting episode can be measured through neural activity, physiology, behavior, and report, but no single channel is a complete ground truth or a universal emotion detector. We distinguish functional-circuit, distributed-representational, predictive-interoceptive, and general-arousal model families, and propose ORIGIN, an organismal reconstruction framework for their falsifiable comparison. The proposed preregistered within-subject study combines fMRI, EEG, pupillometry, electrodermal activity, electrocardiography, respiration, behavior, salivary cortisol, and sparse self-report across threat, reward, social-evaluation, and mild respiratory interoceptive-calibration contexts. Nested held-out prediction, calibration, reliability, and negative controls determine whether an integrated latent-state model explains data beyond stimulus category, task demand, global arousal, and respiratory mechanics. This is a design-only report: it introduces no new participant data or empirical results.
\end{abstract}

\begin{IEEEkeywords}
affective neuroscience, emotion, interoception, allostasis, fMRI, autonomic physiology, predictive modeling, causal inference
\end{IEEEkeywords}

\section{Introduction}

Human emotions organize action under uncertainty, alter perception and learning, regulate internal physiology, and support communication. The everyday question ``Where do emotions originate?'' is therefore scientifically important, but potentially misleading. It suggests that an emotion should have one physical birthplace, analogous to a component in a machine. Decades of lesion, stimulation, neuroimaging, behavioral, autonomic, endocrine, and computational work instead indicate that an emotion episode is assembled across interacting systems. A more precise question is: \emph{which biological processes generate the coordinated neural, bodily, behavioral, and report-level changes observed in an emotion episode, and how can competing mechanistic models be tested?}

This report constrains the answer to natural science. It uses affective and systems neuroscience, autonomic and endocrine physiology, computational modeling, and comparative evidence. It does not use philosophical theories as causal mechanisms, nor does it reduce an episode to a self-report label. The report also avoids a false choice between a brain-only account and a body-only account. Interoceptive inputs, autonomic outputs, immune and endocrine context, and neural prediction are coupled in a control loop. An emotion may be experienced consciously, but conscious report, facial expression, physiology, and circuit activity remain partially dissociable measurements.

The first contribution is an evidence-based reframing. We distinguish evolved functional systems from named human emotion concepts and show why neither an amygdala response nor a heart-rate change is a direct readout of fear, anger, pride, or guilt. The second contribution is a research framework named \origin{}: Organismal Reconstruction of Integrated Generative Neuroaffect. It turns the question of origin into a set of observable state variables, regulated bodily processes, learned predictions, distributed circuit features, and competing null models. The third contribution is a preregistered experiment design intended to compare explanations rather than merely produce another association map.

\section{What Is the Explanandum?}

\subsection{Separate the state, episode, category, and measurement}

The word ``emotion'' names several related but nonidentical targets. Adolphs argues that affective neuroscience must distinguish biological emotion states, conscious experiences, expressions, and concepts \cite{Adolphs2016}. This distinction is not semantic housekeeping; it determines what can be measured and what a causal conclusion means. An organism may deploy defensive autonomic and motor responses without reporting fear, and a person may label a state as fear while physiological activation is low or differently configured. A named category can be learned and communicated while the processes maintaining energy balance and action readiness continue outside awareness.

Table~\ref{tab:targets} defines the four levels used here. Their separation makes it possible to ask which mechanisms are shared across species, which processes vary with context and learning, and which measurements are appropriate at a particular temporal scale. For example, an electrodermal response is informative about sympathetic sweat-gland activity but not sufficient to identify a category. A BOLD pattern can reveal distributed coactivation at seconds-scale resolution but cannot by itself establish that the activated tissue caused an experience. A self-report supplies a privileged record of a participant's accessible appraisal, but reporting also relies on language, attention, memory, and task demand.

\begin{table}[!t]
\caption{Operational targets in an emotion-origin study.}\label{tab:targets}
\centering
\footnotesize
\begin{tabular}{p{0.18\columnwidth}p{0.66\columnwidth}}
\toprule
Target & Operational definition and non-equivalence \\
\midrule
Affective state & Latent, time-varying regulation of valence, arousal, bodily control, and action readiness. It is inferred from multiple channels rather than directly observed. \\
Emotion episode & A context-bound configuration of affect, perception, memory, physiology, action, and sometimes conscious experience. \\
Emotion category & A learned label or concept, such as fear or pride, applied to an episode. It is an outcome and possible predictor, not the entire mechanism. \\
Measurement channel & Neural, autonomic, endocrine, behavioral, or report variable. Each measures a component with its own noise, latency, and construct limits. \\
\bottomrule
\end{tabular}
\end{table}

\subsection{A generative systems question}

Let $\mathbf{z}_{i,t}$ denote an unobserved affective-regulatory state for person $i$ at time $t$. It has dimensions for expected and current bodily condition, sensory context, learned value, uncertainty, and action policy. A minimal state transition model is
\begin{equation}
\mathbf{z}_{i,t}=\mathbf{A}\mathbf{z}_{i,t-1}+\mathbf{B}\mathbf{u}_{i,t}+\mathbf{C}\mathbf{b}_{i}+\boldsymbol{\omega}_{i,t},
\label{eq:state-transition}
\end{equation}
where $\mathbf{u}_{i,t}$ is experimentally controlled input, $\mathbf{b}_{i}$ encodes stable or slowly changing individual properties, $\mathbf{A}$ represents state dynamics, $\mathbf{B}$ and $\mathbf{C}$ are effect matrices, and $\boldsymbol{\omega}_{i,t}$ is process noise. Equation~\eqref{eq:state-transition} is not a claim that the brain implements a simple linear system. It defines a falsifiable modeling target: past state, input, and individual context should jointly improve prospective prediction of current multichannel observations.

Observed data $\mathbf{y}^{(m)}_{i,t}$ for modality $m$ follow a modality-specific observation model,
\begin{equation}
\mathbf{y}^{(m)}_{i,t}=\mathbf{L}^{(m)}\mathbf{z}_{i,t}+\mathbf{D}^{(m)}\mathbf{q}_{i,t}+\boldsymbol{\epsilon}^{(m)}_{i,t},
\label{eq:observation-model}
\end{equation}
where $\mathbf{L}^{(m)}$ maps latent state to a measurement, $\mathbf{q}_{i,t}$ contains nuisance and control variables such as movement, visual energy, motor output, breathing, and task difficulty, $\mathbf{D}^{(m)}$ maps these variables to the measurement, and $\boldsymbol{\epsilon}^{(m)}_{i,t}$ is measurement noise. Equation~\eqref{eq:observation-model} formalizes why one measurement cannot settle the origin question. The same latent state can appear differently across sensors, and the same sensor signal can be caused by non-affective processes.

\section{Evidence for Distributed Biological Origins}

\subsection{Evolved regulatory systems constrain emotion episodes}

Comparative and causal work establishes that mammals possess neural systems for defense, reward pursuit, pain avoidance, attachment-related behavior, arousal, and autonomous regulation. These systems are functionally consequential because perturbing them alters action, learning, physiology, and survival-relevant choices. The amygdala, hypothalamus, periaqueductal gray, ventral striatum, ventral tegmental area, brainstem nuclei, and cortical pathways are central components, with dense recurrent connectivity. Reviews of amygdala anatomy and function emphasize its participation in sensory valuation, associative learning, autonomic and endocrine control, and adaptation to environmental change \cite{Phelps2005,Simic2021}.

However, identifying a circuit that controls a defensive action does not identify a universal ``fear center.'' LeDoux's survival-circuit account separates neural systems that detect and respond to threat from the conscious feeling that a person may call fear \cite{LeDoux2012}. This protects against a common reverse inference: if the amygdala is active, then fear must be present. The correct inference is narrower: a task recruited a region with known participation in functions that can include salience learning, attention, uncertainty, autonomic coordination, and threat-related behavior. Whether fear was experienced, and how it was constructed in context, needs additional evidence.

Reward research gives a parallel lesson. Berridge and Kringelbach distinguish components such as incentive salience, hedonic impact, learning, and motivation \cite{Berridge2008}. Dopamine-related signals can support prediction and wanting without directly equaling pleasure. Orbitofrontal, ventromedial prefrontal, cingulate, amygdalar, striatal, and brainstem interactions connect valued sensory information to goal-directed action \cite{Rolls2023}. The existence of partially specialized mechanisms therefore supports a functional-circuit model, but it does not entail a one-circuit/one-category taxonomy of human feelings.

\subsection{Distributed patterns, not invariant locations}

Meta-analytic evidence challenges strict localization. Reviewing neuroimaging studies, Lindquist et al. found little support for a consistent, specific region dedicated to each discrete emotion category and instead described interacting systems involved in both emotional and non-emotional operations \cite{Lindquist2012}. This conclusion does not mean that neural organization is undifferentiated. It means that a region's contribution depends on network, task, prior learning, bodily state, and temporal phase.

Multivariate studies provide complementary evidence. Kragel and LaBar decoded spontaneous emotional states from distributed brain activity \cite{Kragel2016}; Saarimaki et al. reported that a broad set of induced categories could be classified from widespread cortical and subcortical activity patterns \cite{Saarimaki2018}. These studies support the hypothesis that category-relevant information is distributed. They do not establish that a classifier's features are necessary causes of an experience, because classification can exploit correlates, task structure, sensory features, or report-related processes. The experimental design in Section~\ref{sec:design} therefore makes general arousal, stimulus, motor, and report timing explicit competitors.

\begin{figure*}[!t]
\centering
\fbox{\begin{minipage}{0.94\textwidth}
\centering\small
\textbf{ORIGIN mechanistic schematic.} Exteroceptive cues, context, and memory interact with interoceptive input and allostatic demand. Distributed subcortical, brainstem, insular, cingulate, prefrontal, striatal, and hippocampal systems update a latent regulatory state. The state coordinates autonomic/endocrine change, perception, action selection, learning, and optional report. Neural activity, physiology, behavior, and self-report are treated as distinct observations of the process. A matched arousal/task-demand branch is modeled explicitly rather than explained away after analysis.
\end{minipage}}
\caption{A distributed brain-body account replaces a single-location search with testable state transitions and multichannel observations. The diagram is conceptual and makes no claim that every arrow is a direct monosynaptic pathway.}\label{fig:origin-schematic}
\end{figure*}

\subsection{Interoception, autonomic control, and allostasis}

Every emotion episode occurs in a body. Cardiovascular, respiratory, gastrointestinal, immune, endocrine, and thermoregulatory signals reach the nervous system continuously. Interoception is the central processing of these visceral-afferent signals, while descending autonomic and endocrine pathways help regulate the same body. Schulz characterizes stress and interoception as bidirectional communication along brain-body axes, including the hypothalamic-pituitary-adrenocortical axis and autonomic nervous system \cite{Schulz2015}. Craig highlights the anterior insula as an important part of the pathway from bodily signaling to awareness, while not claiming it operates in isolation \cite{Craig2009}.

The allostasis concept sharpens this biological role. A regulatory system must anticipate energy, temperature, fluid, oxygen, social safety, and other demands before deviation threatens viability. Predictive-interoceptive accounts propose that the brain uses generative models to anticipate bodily causes and update them using interoceptive and exteroceptive prediction errors \cite{Barrett2017,Kleckner2017}. A compact prediction-error term is
\begin{equation}
\boldsymbol{\delta}_{i,t}=\mathbf{s}^{\mathrm{body}}_{i,t}-\widehat{\mathbf{s}}^{\mathrm{body}}_{i,t|t-1},
\label{eq:interoceptive-error}
\end{equation}
where $\mathbf{s}^{\mathrm{body}}_{i,t}$ is a vector of measured or inferred bodily input, $\widehat{\mathbf{s}}^{\mathrm{body}}_{i,t|t-1}$ is the predicted input before the observation, and $\boldsymbol{\delta}_{i,t}$ is the residual. Equation~\eqref{eq:interoceptive-error} becomes scientifically meaningful only when prediction sources, time windows, sensory reliability, and alternative arousal explanations are specified.

Interoception is multidimensional. Accuracy on a cardiac task, confidence in bodily attention, interpretation of breathlessness, and neural responses to heartbeats can dissociate \cite{Suksasilp2022}. Therefore, the present program measures several channels and does not treat heartbeat detection or heart-rate variability as a summary of emotional ability. Cortisol is included only as a slow session-context signal; its temporal dynamics make it inappropriate as a trial-by-trial detector.

\subsection{Learning, memory, and contextual construction}

Context changes the meaning of a signal. A rapid heartbeat may accompany running, anticipation, social evaluation, heat, or threat. Associative learning connects cues to outcomes; hippocampal systems contribute context and memory; prefrontal and striatal pathways modify value and action policy. Constructed-emotion and active-inference accounts suggest that categories and past experience help form predictions that organize otherwise ambiguous sensory and bodily input \cite{Barrett2017}. This is an empirical claim about learning, inference, and biology, not a claim that the body is irrelevant or that arbitrary labels create physiology.

The prediction is conditional. If contextual priors change reports and joint neural-physiological state estimates after sensory intensity, motor requirements, and attention are matched, a predictive-interoceptive model gains support. If only stimulus category and specialized circuit features matter, the functional-circuit model gains support. If both perform no better than a task-demand model, neither emotion-specific account receives support. These discriminations motivate the prospective experiment.

\subsection{Time, scale, and the limits of cross-modal correspondence}

An origin account must respect time. Neural spiking and local field potentials change on millisecond scales; event-related EEG components and pupil responses have different lags; electrodermal responses are delayed and slow; BOLD signals convolve neural and vascular processes over seconds; cortisol changes over minutes. These measures cannot be aligned by placing all of them in one trialwise correlation table. They must be modeled with response functions appropriate to their physiology, as in Equation~\eqref{eq:autonomic-response}. A failure to accommodate time can manufacture apparent disagreement between channels or make a common response look emotion-specific.

Spatial scale also matters. A cell population, a subnucleus, a macroscopic fMRI voxel, and a large-scale functional network are different scientific objects. The amygdala is not homogeneous, nor is the insula, cingulate cortex, or striatum. Conversely, a whole-brain classifier may predict a report while hiding the causal pathway by which a perturbation would change action. The program therefore uses predeclared regions and systems for hypothesis tests, complemented by whole-brain multivariate analyses whose interpretation remains predictive rather than localized.

Finally, a state may be regulated before it is reported. An autonomic change or avoidance decision can precede a category label, and a label can retrospectively reorganize memory of the episode. The appropriate evidence hierarchy moves from measurement reliability, to within-task association, to held-out prediction, to replication, and only then to causal perturbation. No one rung licenses all of the conclusions available at later rungs. This hierarchy is especially important when animal work constrains a circuit's action function but human work asks about conscious, social, or language-mediated episodes.

\section{Competing Models and the ORIGIN Framework}

\subsection{Model families}

Table~\ref{tab:models} converts recurrent theoretical contrasts into testable models. The models are not mutually exclusive at every biological scale. Conserved circuits can implement components inside a distributed predictive system. The purpose of comparison is to determine which components add reliable, generalizable explanatory value for a specified task and measurement set.

\begin{table*}[!t]
\caption{Competing model families and preregistered discriminative predictions.}\label{tab:models}
\centering
\footnotesize
\begin{tabular}{p{0.16\textwidth}p{0.225\textwidth}p{0.225\textwidth}p{0.18\textwidth}}
\toprule
Model family & Mechanistic claim & Key prediction & Critical failure condition \\
\midrule
Functional-circuit & Recurrent challenges recruit partially specialized regulatory/action pathways. & Threat, reward, and social-loss domains retain partly separable pathway and physiological features after common controls. & A domain label adds no generalizable value beyond common arousal and demand features. \\
Distributed-representational & Categories and dimensions are encoded in multiregion multichannel configurations. & Held-out persons and contexts show calibrated decoding using distributed features, not one region. & Decoding is confined to training folds or is explained by sensory/motor/report confounds. \\
Predictive-interoceptive & Contextual priors and expected/actual body state shape the latent state through learning and allostasis. & Cue probability and mild inspiratory load/sham improve joint prediction beyond respiratory mechanics and arousal. & Prior and interoceptive terms do not improve held-out prediction after controls. \\
General-arousal/task-demand null & Apparent signatures reflect salience, attention, effort, novelty, or response preparation. & Matched non-affective controls account for observed multichannel changes. & Condition-specific, replicable multimodal structure remains after all control variables. \\
\bottomrule
\end{tabular}
\end{table*}

\subsection{ORIGIN principles}

The ORIGIN framework has six principles. \textbf{O} is observable operationalization: the study separately defines neural, physiological, behavioral, and report outcomes. \textbf{R} is regulation and allostasis: body control is represented as a mechanism rather than automatically regressed away. \textbf{I} is interoception and neuromodulation: sensing and regulating the body are measured through channels with different latencies. \textbf{G} is generative prediction and learning: task cues manipulate priors and expected outcomes. The second \textbf{I} is integrated distributed circuits: statistical features span predeclared networks and whole-brain tests. \textbf{N} is null and competing models: general arousal, attention, motor preparation, and demand characteristics receive explicit models.

This framework also specifies what it cannot show. A strong held-out model does not prove that its variables constitute a complete theory of subjective experience. A multimodal association cannot substitute for causal perturbation. A group-level pattern cannot diagnose an individual. These limits are not generic caveats; they determine the analysis and translation rules in Sections~\ref{sec:design} and \ref{sec:limits}.

\subsection{Causal evidence hierarchy}

The program organizes claims by the evidence they require. First, sensor validity asks whether an instrument is stable, temporally interpretable, and responsive to the stated manipulation. Second, observational association asks whether a predeclared neural or bodily feature covaries with task context while known confounds are modeled. Third, held-out prediction asks whether a model calibrated on one subset anticipates a new person, trial, or context. Fourth, replication asks whether the effect survives an independent sample, apparatus, and analysis team. Fifth, causal evidence asks whether an ethically justified intervention changes a registered outcome with an explicit identification strategy.

The hierarchy has a practical consequence. A high cross-validated decoding score may be valuable measurement science, but it is not a license to call a selected voxel the origin of a feeling. Similarly, stimulation can provide a causal result for an action, autonomic output, or report in a narrow context without proving that the stimulated node is the sole source of the full episode. The proposed study is intentionally located at the first three levels. It creates a compact, shareable assay that could support later clinical or perturbational studies only if it generalizes.

\section{Preregistered Multimodal Experiment}\label{sec:design}

\subsection{Objectives and hypotheses}

The primary objective is to determine whether a joint latent-state model improves held-out prediction of brain, body, action, and report data over a general-arousal/task-demand null. The study is explicitly \designonly{}: no participants have been enrolled and no effect size, classifier accuracy, biomarker, or treatment result is claimed. The study enrolls 180 healthy adults aged 18 to 40. The confirmatory analysis set requires at least 150 participants with at least 80\% valid trials; if the threshold is missed, results are reported as under-target and recruitment is not expanded after outcome inspection. Exclusions, missingness handling, and the complete analysis code are fixed before enrollment.

The primary hypothesis is that distributed predictive-interoceptive state features provide a prespecified improvement in expected log predictive density over the null. Secondary hypotheses test whether functional-circuit features retain domain-specific predictive value and whether the models generalize across held-out participants and contexts. An informative null result occurs if the emotion-oriented models fail to beat a matched stimulus/arousal model, or if apparent performance disappears under context-general validation.

\subsection{Tasks and experimental controls}

Participants complete four counterbalanced domains: threat uncertainty, reward approach, social evaluation/loss, and interoceptive calibration. Threat trials use explicitly low-intensity aversive versus neutral sounds with probabilistic cues. Reward trials use small monetary gains, omissions, and neutral outcomes with identical timing and response demands. Social-evaluation/loss trials use standardized non-personal feedback and symbolic loss cues. The health-screened interoceptive module presents mild inspiratory resistance or sham following learned probability cues, while airway pressure and respiratory measures verify actual bodily input. Each domain contains learned cue contingencies and a preregistered reversal phase, enabling tests of context and prediction error.

The design includes non-affective controls matched for visual complexity, motor response frequency, working-memory load, novelty, and expected value where possible. The interoceptive module crosses cue probability with actual load/sham delivery and includes respiratory-mechanics covariates. Ratings of valence, arousal, uncertainty, bodily intensity, and an optional category label are sampled on sparse randomized trials. Sparse sampling minimizes the chance that continuous verbalization changes state dynamics. The label is modeled as one observation in Equation~\eqref{eq:observation-model}, not as the definition of the state.

\begin{table}[!t]
\caption{Prospective data-acquisition plan.}\label{tab:acquisition}
\centering
\footnotesize
\begin{tabular}{p{0.20\columnwidth}p{0.65\columnwidth}}
\toprule
Channel & Purpose and limitation \\
\midrule
BOLD fMRI & Estimates distributed systems including cortical and subcortical structures at a hemodynamic timescale; correlational and motion-sensitive. \\
MRI-compatible EEG & Captures event-locked cortical timing; requires explicit gradient and pulse-artifact quality control. \\
Pupil, EDA, ECG, respiration, pressure & Samples autonomic, orienting, respiratory, and delivered-load dynamics; non-specific and analyzed with modality-specific lags. \\
Action and response time & Captures choice, avoidance/approach, and motor preparation; can be a confound and an outcome. \\
Salivary cortisol & Baseline and post-session context measure only; not a trial-level emotion label. \\
Sparse report & Records accessible appraisal and category use; affected by language, attention, memory, and demand. \\
\bottomrule
\end{tabular}
\end{table}

\subsection{Locked multimodal state estimation}

The primary analysis uses a hierarchical state-space model. The latent state has six locked dimensions: expected value, uncertainty, interoceptive prediction error, regulatory arousal, context/memory state, and approach-avoidance policy. The functional-circuit model uses prespecified amygdala, hypothalamus, periaqueductal gray, ventral striatum, anterior insula, dorsal anterior cingulate, ventromedial prefrontal, and hippocampal features. The distributed model uses a fixed cortical-subcortical parcel set, and the null model uses visual/auditory energy, reaction time, motor response, task difficulty, novelty, and measured respiratory mechanics. Individual parameters partially pool toward group distributions while retaining person-specific uncertainty. For a scalar endpoint $r_{i,t}$, such as a report dimension or a predeclared physiological feature, the hierarchical component can be written as
\begin{equation}
r_{i,t}=\alpha_i+\boldsymbol{\beta}^{\mathsf{T}}\mathbf{z}_{i,t}+\boldsymbol{\gamma}^{\mathsf{T}}\mathbf{q}_{i,t}+\eta_{i,t},\quad \alpha_i\sim\mathcal{N}(\mu_{\alpha},\sigma_{\alpha}^{2}),
\label{eq:hierarchical-endpoint}
\end{equation}
where $\alpha_i$ is a participant intercept, $\boldsymbol{\beta}$ links latent state to the endpoint, $\boldsymbol{\gamma}$ adjusts for prespecified confounds, and $\eta_{i,t}$ is residual noise. Equation~\eqref{eq:hierarchical-endpoint} prevents an average effect from being mistaken for an invariant personal signature.

Modality-specific temporal response functions are estimated before fusion. For example, an event-locked autonomic response is represented as
\begin{equation}
a_{i}(\tau)=\sum_{k=1}^{K}\theta_{i,k}h_k(\tau)+e_i(\tau),
\label{eq:autonomic-response}
\end{equation}
where $a_i(\tau)$ is an autonomic waveform relative to an event, $h_k(\tau)$ are preregistered basis functions at latency $\tau$, $\theta_{i,k}$ are participant coefficients, and $e_i(\tau)$ is error. Equation~\eqref{eq:autonomic-response} avoids forcing electrodermal, cardiac, pupil, EEG, and BOLD signals into a common instantaneous time bin.

Data quality is an output, not an afterthought. The protocol records fMRI framewise displacement and censoring, MRI distortion and physiological-regressor decisions, EEG rejected channels/trials and residual artifacts, eye-tracking loss, ECG detection failures, respiration quality, assay timing, and reasons for missing data. Feature extraction, scaling, imputation, and selection are learned only on training partitions. The released analysis will include complete-case and prespecified missingness sensitivity analyses.

\subsection{Prediction, calibration, and reliability}

For each model $M$, the principal held-out score is expected log predictive density,
\begin{equation}
\mathrm{ELPD}(M)=\sum_{(i,t)\in\mathcal{T}_{\mathrm{test}}}\log p_M\!\left(\mathbf{y}_{i,t}\mid\mathcal{D}_{\mathrm{train}}\right),
\label{eq:elpd}
\end{equation}
where $\mathcal{T}_{\mathrm{test}}$ contains held-out participant-context observations and $\mathcal{D}_{\mathrm{train}}$ excludes them. Equation~\eqref{eq:elpd} rewards calibrated probability assigned to unseen data rather than a descriptive fit to the same observations used for training. The registered primary comparison averages equally weighted fMRI, EEG, autonomic, behavior, and report outcome-family ELPDs. ORIGIN support requires $\Delta\mathrm{ELPD}\geq 0.02$ nats per trial, a bootstrap 95\% lower confidence bound above zero, and no Brier-score worsening larger than 0.005 relative to the null.

For categorical or binned predictions, calibration is assessed by the Brier score,
\begin{equation}
\mathrm{BS}=\frac{1}{N}\sum_{n=1}^{N}\sum_{c=1}^{C}\left(\widehat{p}_{n,c}-\mathbb{I}[y_n=c]\right)^2,
\label{eq:brier}
\end{equation}
where $N$ is the number of held-out observations, $C$ is the number of outcome classes, $\widehat{p}_{n,c}$ is predicted probability, and $\mathbb{I}[y_n=c]$ is an indicator. Lower values in Equation~\eqref{eq:brier} indicate more accurate probability forecasts. Test-retest and split-half stability are summarized with an intraclass correlation coefficient,
\begin{equation}
\mathrm{ICC}=\frac{\sigma^{2}_{\mathrm{between}}}{\sigma^{2}_{\mathrm{between}}+\sigma^{2}_{\mathrm{within}}},
\label{eq:icc}
\end{equation}
where the numerator is between-person variance and the denominator adds within-person error variance. Equation~\eqref{eq:icc} is reported for measurement components and individual-level model outputs; unreliability limits interpretation even if a group mean differs from zero.

\subsection{Model comparison and falsification}

Cross-validation is nested. Outer folds hold out both people and a subset of contexts, while inner folds select tuning parameters and any feature reduction. The model set is locked before outcome analysis: functional-circuit, distributed-representational, predictive-interoceptive, and general-arousal/task-demand null. Negative controls include non-affective high-effort blocks, permutation of condition labels within allowable strata, and classifiers restricted to visual, motor, and respiratory-mechanics nuisance features. A result is not called emotion-specific if these controls perform equivalently.

Before recruitment, the project will release the locked versioned preregistration containing hypotheses, the sampling rule, task randomization, hardware settings, preprocessing pipeline, all model equations, outcome-family weighting, the $0.02$-nat margin, the Brier guardrail, and the exact rules for deviations. Deviations are not forbidden; they are recorded with rationale and analyzed separately from the confirmatory path. A synthetic pilot dataset will test whether the pipeline can recover known simulated effects and whether missing channels or trial loss produce unstable estimates. This is a software and design validation, not evidence about human emotion.

The analysis reports both aggregate and individual uncertainty. Group estimates summarize regularities but cannot imply that every participant expressed the same autonomic or neural trajectory. Participant-level posterior predictive checks examine whether the model reproduces each person's timing and covariance structure. Cross-context validation is decisive: training on threat and testing on a new threat cue probes within-domain robustness, while training across threat and reward and testing social evaluation probes whether a proposed latent variable has any broader explanatory range. These tests convert vague claims of ``general emotion'' into a measurable scope statement.

Figure~\ref{fig:inference-pipeline} shows the evidential sequence. The primary study makes no causal claim from ordinary fMRI or physiology. A future perturbation study should identify a concrete intervention $A$, a target outcome $Y$, and a population before estimating
\begin{equation}
\tau=\mathbb{E}\left[Y\left(A=1\right)-Y\left(A=0\right)\right],
\label{eq:causal-estimand}
\end{equation}
where $Y(A=1)$ and $Y(A=0)$ are counterfactual outcomes under intervention and control. Equation~\eqref{eq:causal-estimand} requires exchangeability, consistency, positivity, safety, and an intervention that can be ethically delivered. None of these assumptions is satisfied merely by observing BOLD activation.

\begin{figure*}[!t]
\centering
\fbox{\begin{minipage}{0.94\textwidth}
\centering\small
\textbf{Prospective inference pipeline.} Preregister tasks, sensor quality rules, model set, and controls $\rightarrow$ acquire synchronized brain-body-action-report data $\rightarrow$ preprocess within training folds $\rightarrow$ fit modality-specific and joint models $\rightarrow$ test held-out persons and contexts $\rightarrow$ assess ELPD, calibration, reliability, and negative controls $\rightarrow$ classify evidence as support, inconclusive, or contradiction for each model family. Causal claims require a later intervention module with explicit assumptions.
\end{minipage}}
\caption{The study's decision flow. It protects against data leakage, reverse inference, and post hoc theory selection.}\label{fig:inference-pipeline}
\end{figure*}

\section{Ethics, Interpretation, and Translation}\label{sec:limits}

\subsection{Participant welfare and data governance}

The proposed task uses low-risk stimuli, participant-controlled pauses, discontinuation without penalty, content warnings, trained monitoring, and a debrief. The respiratory module uses only mild, health-screened inspiratory resistance with real-time participant stop control and airway-pressure monitoring. The protocol excludes severe stress induction, deceptive threat, pharmacologic challenge, and invasive stimulation. These boundaries are part of measurement validity: an ethically disproportionate manipulation would change the scientific target and reduce generalizability to everyday emotion regulation.

Neural and physiological recordings are sensitive data. Data governance should use coded identifiers, minimal access, separation of direct identifiers from research data, auditable access controls, and consent language that distinguishes aggregate research use from clinical use. The study should not produce an individual emotion score for employment, education, insurance, policing, or clinical triage. Population-level prediction does not automatically support individual-level inference.

\subsection{Interpretation guardrails}

Table~\ref{tab:guardrails} gives precommitted language rules. They matter because the public receives claims such as ``the amygdala causes fear'' or ``a wearable detects emotion'' more readily than the uncertainty and scope conditions that make those claims defensible. The present report uses evidence to constrain mechanisms, not to sell a detector.

\begin{table*}[!t]
\caption{Interpretation guardrails for evidence and anticipated outcomes.}\label{tab:guardrails}
\centering
\footnotesize
\begin{tabular}{p{0.23\textwidth}p{0.33\textwidth}p{0.35\textwidth}}
\toprule
Observation & Permitted interpretation & Prohibited interpretation \\
\midrule
Amygdala or insula association & A task recruited a region participating in functions relevant to salience, interoception, learning, or regulation. & A single region proves a specific conscious emotion. \\
Distributed held-out decoding & A multichannel pattern contains generalizable information under this task and population. & The pattern is a universal biological fingerprint or an individual diagnostic. \\
Autonomic change & The body changed in temporal relation to the event and contributes useful state information. & Heart rate, EDA, HRV, pupil, or cortisol uniquely identifies an emotion category. \\
Prior/reversal model improvement & Contextual learning contributes predictive value beyond the registered controls. & Emotions are arbitrary or entirely created by language. \\
Causal perturbation result & The intervention changed an outcome in the defined population under stated assumptions. & A manipulation identifies the sole origin of an emotion. \\
\bottomrule
\end{tabular}
\end{table*}

\subsection{Scientific and translational limits}

The study tests a limited task battery in healthy adults. Threat uncertainty, reward approach, and social evaluation/loss are useful laboratory probes, but they do not exhaust human emotion. The sample cannot establish developmental mechanisms, sex-linked physiology, cross-cultural variation, psychiatric pathophysiology, or cross-species equivalence without dedicated designs. Simultaneous EEG-fMRI improves temporal complementarity but creates artifact and analytic burdens. Salivary cortisol is temporally slow and intentionally cannot adjudicate trial-level models.

The proposal also faces the reproducibility problem. Small samples, flexible preprocessing, and post hoc model selection can overstate brain-behavior relations \cite{Button2013}. The response is not to abandon complex measurement; it is to preregister, validate out of sample, quantify uncertainty, publish exclusions, test negative controls, and plan independent replication. A model that cannot generalize to a new context should be described as task-bound, not as an origin theory.

If the integrated approach succeeds, its near-term value is mechanistic measurement science. It could help refine hypotheses about dysregulated brain-body control in anxiety, depression, trauma-related syndromes, chronic pain, or autonomic disease. It must not be treated as a clinical biomarker until an independent clinical program demonstrates reliability, calibration, equity, utility, and safety for the intended use. The distinction between a scientific association and a medical decision is fundamental.

\section{Research Agenda for Resolving the Origin Question}

\subsection{From descriptive maps to intervention-ready mechanisms}

An emotion-origin program should progress through distinct scientific milestones. The first milestone is measurement validity. Researchers need to demonstrate that each neural, autonomic, endocrine, behavioral, and report feature is reliably acquired, is temporally interpretable, and is not dominated by simple artifacts. The second is representational validity: a model must anticipate observations in held-out data and must disclose where it generalizes and where it fails. The third is mechanistic validity: a model should correctly predict changes produced by a specific perturbation. This order prevents a familiar error in neuroscience, where a visually compelling activation map is promoted to a cause before its measurement reliability or out-of-sample predictability are known.

\subsection{Cross-species translation with explicit limits}

An effective bridge uses homologous task components instead of assuming homologous folk categories. Examples include probabilistic threat learning, approach under reward uncertainty, avoidance action, reversal learning, social-safety signaling, and interoceptive perturbation. The bridge then asks whether neural dynamics, physiological control, and behavior follow comparable computational variables. This approach respects shared mammalian biology while avoiding anthropomorphic overreach. It also provides a route for reciprocal validation: animal perturbation should predict an observable feature in the human assay, and human contexts that invalidate the feature should update the animal model.

\subsection{Development, plasticity, and heterogeneity}

No origin theory is complete without development. Prenatal influences, maturation of autonomic and cortical systems, attachment-relevant caregiving, stress exposure, sleep, nutrition, learning, and environmental uncertainty can alter the gain and calibration of brain-body regulation. A systems account predicts heterogeneity rather than a uniform circuitry-to-category map. Two people can both report fear while one exhibits strong cardiovascular mobilization and another exhibits freezing or dissociative responses; neither pattern should be declared more genuine solely because it resembles a group mean.

\subsection{Open and adversarial science}

The field benefits from pre-registered adversarial collaboration. Investigators representing functional-circuit, distributed-representational, predictive-interoceptive, and general-arousal accounts should agree in advance on tasks, control conditions, feature families, evaluation scores, and interpretations of possible outcomes. This is not a demand for theoretical consensus. It is a way to ensure that a result unfavorable to one account is genuinely diagnostic instead of dismissed as an unanticipated design choice.

To make such collaboration practical, the study should release a synthetic dataset before collection, task code, randomization seeds, hardware timing tests, preprocessing containers, and a fully specified analysis notebook. After primary analysis, deidentified summary data and feature-level derivatives should be shared subject to consent, governance, and re-identification risk. Independent teams should reproduce the analysis from raw or controlled-access data and then repeat the task with a new scanner site or population. Replication must target both average effects and reliability: a nominally significant group contrast is inadequate if participant-level state estimates are unstable.

\subsection{Decision rules for progress}

The program advances only when evidence passes all relevant gates. A candidate emotion-state feature must be technically reliable, add held-out predictive value beyond a matched control model, remain calibrated under a context shift, and replicate. A causal claim must additionally identify and validate an intervention pathway. A translational claim must show clinical utility, not merely association with a diagnosis. This gate sequence is intentionally demanding because emotion-related measures are attractive targets for surveillance, manipulation, and overconfident clinical interpretation.

\section{Discussion}

The evidence supports a layered answer to the origin question. At an evolutionary level, recurrent survival and reproductive challenges have shaped circuits for defense, reward, pain avoidance, arousal, caregiving, and social regulation. At an organismal level, these circuits are embedded in loops that sense and regulate internal physiology. At a systems level, cortical, subcortical, brainstem, hippocampal, and neuromodulatory processes integrate cues, value, action possibilities, and memory. At a computational level, learning and prediction help resolve ambiguity between the state of the body and the likely causes in the world. At a report level, attention, language, category knowledge, and social communication make part of this process accessible and nameable.

This layered account explains why a single answer such as ``the amygdala,'' ``hormones,'' or ``the body'' is incomplete. It also resists the opposite error of treating distributed processing as a non-mechanistic cloud. Distribution is testable when state variables, pathways, timescales, inputs, and falsification criteria are specified. The ORIGIN proposal makes those specifications explicit through Equations~\eqref{eq:state-transition}-\eqref{eq:causal-estimand}, the model families in Table~\ref{tab:models}, and the registered controls.

There is a genuine scientific dispute about how much explanatory work is performed by specialized functional systems, distributed representations, and predictive-interoceptive inference. The appropriate response is neither to declare a winner from a single imaging study nor to frame the dispute as a purely conceptual preference. The models make different predictions about context reversals, bodily prediction, cross-context generalization, and the residual explanatory value of specialized pathway features. The proposed design measures those predictions with multiple channels and requires the emotional account to outperform a general arousal explanation.

\section{Conclusion}

Human emotions do not originate in a single brain center or a single chemical. They emerge from distributed, evolving brain-body systems that coordinate regulation, interoception, sensation, learning, memory, valuation, action, and report. Conserved functional circuits provide biologically important constraints, while context and prediction help shape the episode observed in a particular person and moment. The most useful next step is a preregistered multimodal comparison of competing models with held-out prediction and clear causal boundaries. The ORIGIN framework provides that step without converting a complex biological process into an unsupported emotion detector.

\begin{thebibliography}{00}

\bibitem{Adolphs2016}
R. Adolphs, ``How should neuroscience study emotions? By distinguishing emotion states, concepts, and experiences,'' \emph{Social Cognitive and Affective Neuroscience}, vol. 12, no. 1, pp. 24-31, 2017, doi: 10.1093/scan/nsw153.

\bibitem{Phelps2005}
E. A. Phelps and J. E. LeDoux, ``Contributions of the amygdala to emotion processing: From animal models to human behavior,'' \emph{Neuron}, vol. 48, no. 2, pp. 175-187, 2005, doi: 10.1016/j.neuron.2005.09.025.

\bibitem{Simic2021}
G. Simic, M. Tkalcic, V. Vukic, D. Mulc, E. Spanic, M. Sagud, F. E. Olucha-Bordonau, M. Vuksic, and P. R. Hof, ``Understanding emotions: Origins and roles of the amygdala,'' \emph{Biomolecules}, vol. 11, no. 6, Art. no. 823, 2021, doi: 10.3390/biom11060823.

\bibitem{LeDoux2012}
J. E. LeDoux, ``Rethinking the emotional brain,'' \emph{Neuron}, vol. 73, no. 4, pp. 653-676, 2012, doi: 10.1016/j.neuron.2012.02.004.

\bibitem{Berridge2008}
K. C. Berridge and M. L. Kringelbach, ``Affective neuroscience of pleasure: Reward in humans and animals,'' \emph{Psychopharmacology}, vol. 199, no. 3, pp. 457-480, 2008, doi: 10.1007/s00213-008-1099-6.

\bibitem{Rolls2023}
E. T. Rolls, ``Emotion, motivation, decision-making, the orbitofrontal cortex, anterior cingulate cortex, and the amygdala,'' \emph{Brain Structure and Function}, vol. 228, no. 5, pp. 1201-1257, 2023, doi: 10.1007/s00429-023-02644-9.

\bibitem{Lindquist2012}
K. A. Lindquist, T. D. Wager, H. Kober, E. Bliss-Moreau, and L. F. Barrett, ``The brain basis of emotion: A meta-analytic review,'' \emph{Behavioral and Brain Sciences}, vol. 35, no. 3, pp. 121-143, 2012, doi: 10.1017/S0140525X11000446.

\bibitem{Kragel2016}
P. A. Kragel and K. S. LaBar, ``Decoding spontaneous emotional states in the human brain,'' \emph{PLoS Biology}, vol. 14, no. 9, Art. no. e2000106, 2016, doi: 10.1371/journal.pbio.2000106.

\bibitem{Saarimaki2018}
H. Saarimaki, L. F. Ejtehadian, E. Glerean, I. P. J. Jaskelainen, P. Vuilleumier, M. Sams, and L. Nummenmaa, ``Distributed affective space represents multiple emotion categories across the human brain,'' \emph{Social Cognitive and Affective Neuroscience}, vol. 13, no. 5, pp. 471-482, 2018, doi: 10.1093/scan/nsy018.

\bibitem{Schulz2015}
A. Schulz, ``Interoception and stress,'' \emph{Frontiers in Psychology}, vol. 6, Art. no. 993, 2015, doi: 10.3389/fpsyg.2015.00993.

\bibitem{Craig2009}
A. D. Craig, ``How do you feel - now? The anterior insula and human awareness,'' \emph{Nature Reviews Neuroscience}, vol. 10, no. 1, pp. 59-70, 2009, doi: 10.1038/nrn2555.

\bibitem{Barrett2017}
L. F. Barrett, ``The theory of constructed emotion: An active inference account of interoception and categorization,'' \emph{Social Cognitive and Affective Neuroscience}, vol. 12, no. 1, pp. 1-23, 2017, doi: 10.1093/scan/nsw154.

\bibitem{Kleckner2017}
I. R. Kleckner \emph{et al.}, ``Evidence for a large-scale brain system supporting allostasis and interoception in humans,'' \emph{bioRxiv}, 2017, doi: 10.1101/098970.

\bibitem{Suksasilp2022}
C. Suksasilp and S. N. Garfinkel, ``Towards a comprehensive assessment of interoception in a multi-dimensional framework,'' \emph{Biological Psychology}, vol. 168, Art. no. 108262, 2022, doi: 10.1016/j.biopsycho.2022.108262.

\bibitem{Button2013}
K. S. Button \emph{et al.}, ``Power failure: Why small sample size undermines the reliability of neuroscience,'' \emph{Nature Reviews Neuroscience}, vol. 14, no. 5, pp. 365-376, 2013, doi: 10.1038/nrn3475.

\end{thebibliography}

\end{document}
